% LSPV-4 REPORT::LIVEPREFIX1::2026-10-13T22:52:00Z
% !TEX program = xelatex
\documentclass[11pt]{article}
\usepackage[T1]{fontenc}
\usepackage{lmodern,amsmath,amssymb,amsthm,booktabs,xurl,hyperref}
\usepackage[margin=0.9in]{geometry}
\hypersetup{hidelinks}
\setlength{\emergencystretch}{3em}
\newtheorem{theorem}{Theorem}
\newtheorem{lemma}[theorem]{Lemma}
\newcommand{\poly}{\operatorname{poly}}
\newcommand{\SAT}{\operatorname{SAT}}
\title{LIVEPREFIX1: A Charged Post-Disclosure Solver\\
\large Active hash constraints, causal commitments, and a two-tree witness test}
\author{UNIVERSAL\_CORE\_BRIDGE / LSPV-4}
\date{14 October 2026, 07:52 KST}
\begin{document}\maketitle
\paragraph{Status.}
\textbf{CANDIDATE with complete elementary proofs; independent review pending.}
The principal implication is \textbf{CONDITIONAL} on a polynomial decider for
one fixed NP-complete language. The result extends the current wrapper audit
to hashes that participate in actual witness constraints, not just record tags.
It does not establish a contradiction: an impossible advance commitment is
not implied by post-disclosure witness search. \textbf{P versus NP remains OPEN.}

\section{The selected contradiction attempt}
Fix $L=L_{\mathrm{mono22}}$: satisfiable encodings of nonempty proper,
distinct-clause, clause-monotone exact-$(2,2)$ 3-CNF. Malformed encodings reject.
Each clause is entirely positive or entirely negative; the formula need not
have only one polarity. This fixed language is NP-complete \cite{docker}.
Assume a total deterministic uniform decider $A$ with
\[
 T_A(y)\le C_A(|y|+2)^k                                  \tag{1}
\]
for arbitrary fixed $C_A,k$. The intended contradiction was to use a charged
WRAPPER to turn $A$ into a strategy violating the fresh-challenge boundary.
No exponential numerical lower bound or broader-input transfer is required
for a direct contradiction to (1).

The question here is whether making the hash part of the accepted witness
relation, rather than a mere metadata field, supplies that contradiction.
The answer is negative for the explicit bounded local-verifier interface
below. This is not an impossibility theorem for every wrapper design.

\section{Exact active-constraint interface}
Let $N\ge2$ be the original explicit input length. A fixed controller has at
most $N^a$ rounds. At round $i$, it receives a currently disclosed challenge
$r_i$ and a retained history $h_i$. A fixed polynomial-time predicate
\[
 V(x,h_i,r_i,w)=1                                       \tag{2}
\]
checks a response $w$ of padded length $m(N)\le N^b$. In the original
completion-first ordering, this is the next response after the preceding
sample; the sample is not moved before the event that produced it. The full prefix-query
record, including $x,h_i,r_i$ and a tentative prefix of $w$, has length
polynomial in $N$. Lengths and clocks are explicit, not succinct exponential
bounds. The controller, sample acquisition, exact record output, any replay,
and hash updates together have an independently established $O(N^g)$ cost.
Exogenous waits are excluded from active work; a classical elapsed-time
claim additionally needs those waits polynomially bounded.

The verifier may enforce source satisfaction, old commitments, new pinning
constraints, shared witnesses, $H(h_i,w)=r_i$, or another fixed
polynomially checkable hash equation. No assumption about finding hash
preimages or avoiding collisions is used. A digest is not an injective record.
Actual sampled blocks needed for replay are retained as input data.

For a prefix $s$, define the single language
\[
 E_V=\{(x,h,r,s):\exists w\in\{0,1\}^{m(N)},\ s\preceq w,
                         \ V(x,h,r,w)=1\}.             \tag{3}
\]
It lies in NP. Fix a polynomial many-one reduction $\rho:E_V\to L$.
On records generated here, write its charged construction cost as $O(N^p)$
and its output length as $O(N^q)$. These constants depend on the fixed
verifier and encodings, not on $A$ or its exponent. This is an ordinary
NP-completeness use; no local-indexing or near-linear reduction is claimed.
The finite experiments do not implement this monotone compiler.

\begin{theorem}[Post-disclosure prefix selector]\label{thm:prefix}
Under (1), one deterministic uniform algorithm $B_A$ returns the
lexicographically first response satisfying (2), or correctly returns
\textnormal{NO\_EXTENSION}, on every currently supplied round. For the whole
bounded interaction its active time is
\[
 O_A\!\left(N^{a+b+p}+N^{a+b+qk}+N^g\right).             \tag{4}
\]
The algorithm is causal: it uses no later $r_j$. Polynomial arbitrary
workspace is allowed and charged. The claim holds even if the current
challenge and its constraints depend on all prior actual hash values.
\end{theorem}
\begin{proof}
Ask $A(\rho(x,h,r,\epsilon))$. If the answer is NO, no feasible response
exists. Otherwise maintain a prefix $s$ known to extend to a feasible $w$.
Ask whether $s0$ extends. On YES, append zero; on NO, append one, which must
extend by the invariant. After $m(N)$ bits, verify (2) directly and output
$w$. Thus at most $m(N)+1$ decisions are needed per round. Apply (1) to the
$O(N^q)$-bit queries and charge each $O(N^p)$ construction. Summing over
$N^a$ rounds and adding controller work proves (4). All inputs to the
current computation are already disclosed, proving causality.

Constants in (4) may depend on $A,k$ and the fixed reduction, but not $N$.
The same finite program serves all input lengths; it does not choose a
new program or advice for each length. Knowing an optimal $k$ is unnecessary.
\end{proof}

A conservative joint-space ledger is
\[
 S_{B_A}(N)\le S_{\rm controller}(N)+S_\rho(N)+m(N)
                        +S_A(O(N^q)).                  \tag{5}
\]
Retained input, pins, transcript, sample records, reduction buffers and
response output are included. Since space is bounded by time for the
ordinary simulation, (5) is polynomial under (1), not subpolynomial.
No cited small-space SAT exclusion contradicts (4)--(5) \cite{mw}.

\paragraph{What full online success additionally requires.}
The theorem does not assert that a feasible response always exists, or that
an arbitrary locally valid response preserves all future choices. If every
reachable locally valid history and every allowed challenge have a feasible
successor, and terminal local validity is the desired success condition,
induction gives a complete successful interaction. That is an explicit
hereditary-totality hypothesis, not something silently inferred from NP.
If it fails, NO\_EXTENSION identifies an infeasible current response problem;
it is not necessarily UNSAT of the original source.

\paragraph{Decision stages and reference execution.}
For the original rule $e_i=A(F_i)$, sample $r_i$ after completion, then set
$F_{i+1}=C(F_i,H(h_i,e_i,r_i))$, polynomial round count and query size already
give a polynomial own-run simulation under (1). Theorem~\ref{thm:prefix}
adds witness constraints to that closure. It does not reproduce another
implementation's completion times, internal states, sampled blocks, or hash.
Producing a specified reference transcript is a different interface and
requires its actual external observations and charged simulation.

\section{A source-matched two-tree repair}
Retain TWOTREEBRIDGE1's exact construction: split source-graph occurrences;
link consecutive copies with a diamond having edges $xa,xb,ab,ay,by$;
partition them as $\{xa,ab,by\}$ and $\{xb,ay\}$; place source edges in the
second forest; connect each forest's components by private two-edge hub paths.
The current source code is
\path{experiments/review/two_tree_bridge_audit.py}, blob
\path{861b509c25fb2067c59de1695a0d953f31a55f5c}.
This run implements and tests that specified construction, not a replacement
for its pending general independent review.

The target encoding uses the existing eight-bit header, unary universe size
$h$, two membership vectors and two adjacency triangles. Its exact length is
$9+h^2+2h$ bits. The decoder checks lengths, endpoints, tree connectivity and
coverage of the universe. Wrong syntactic records reject; a valid but
incompatible pair remains a legitimate NO instance.

Let $V_I(w)$ check a complete common 3-colouring of the two supplied trees.
This relation is polynomial-time checkable. Its prefix language has (3)'s
form. Therefore (1) supplies a polynomial post-disclosure selector for
arbitrary currently imposed polynomially checkable pins and hash conditions,
not merely the already solved individual trees. This does not create an
unconditional algorithm for the NP-complete compatibility problem.

\begin{lemma}[One revealed colour can be repaired from a witness]\label{lem:palette}
For any compatible pair of trees, any complete common colouring $w$, vertex
$v$, and newly disclosed colour $c\in\{0,1,2\}$, a complete colouring $w'$
with $w'(v)=c$ is computable in $O(h)$ colour operations from $w$.
\end{lemma}
\begin{proof}
Apply the same permutation of the three colour names at every vertex,
swapping $w(v)$ and $c$. Equality and inequality of colours are preserved,
so every edge remains properly coloured. Output and verification have
polynomial encoded cost. No source-solving call is needed.
\end{proof}
This lemma permits global recolouring and has no other immutable colour
pins. It is not a theorem that arbitrary multiple pins can be repaired, or
that previously committed colour names may be silently changed.

If the wrapper instead requires $w(v)$ to be committed before $c$ arrives
and forbids changes, a uniform independent $c$ is matched with probability
exactly $1/3$ by any committed colour, even with unlimited computation.
Both statements apply to the same elementary source relation. Hashing the
old commitment does not turn the latter obligation into ordinary static
colourability. A polynomial decider is required to decide the given problem,
not to predict the later pin.

\section{The load-bearing quantifier distinction}
For a two-round acceptance relation $W(r_0,w_0,r_1,w_1)$, distinguish
\[
 \underbrace{\forall r_0,r_1\ \exists w_0,w_1\ W}_{\text{offline completion}}
 \qquad\text{and}\qquad
 \underbrace{\forall r_0\ \exists w_0\ \forall r_1\ \exists w_1\ W}
 _{\text{causal winning responses}}.                    \tag{6}
\]
The predicate $W\equiv(w_0=r_1)$ makes the first true and the second false.
Static completion can be easy while advance prediction is impossible.
A reduction cannot exchange these orders without a proved source-specific
argument. The example is not evidence against P=NP.

All $2^{16}=65,536$ Boolean relations on four bits were checked. Exactly
$15^4=50,625$ satisfy offline completion. For each fixed $r_0$, the two
candidate $w_0$ rows each admit all future challenges in 9 of their 16
possible four-bit relations. Thus the number of winning eight-bit blocks
is $16^2-7^2=207$, giving $207^2=42,849$ causal relations. Exactly
$7,776$ relations satisfy the offline condition without the causal one.
The exhaustive code agrees with independent direct strategy enumeration.

CONTGAME1 already records the fixed-versus-growing continuation distinction:
a bounded number of future quantifier blocks lies at a fixed PH level;
polynomially many alternating rounds have a PSPACE upper bound, matching
the standard alternation theorem \cite{cks}. The present derivation does not
establish a general PSPACE collapse from P=NP. Requiring a global winning
strategy needs a different reduction, not another name for (3).

\section{What survives as a possible contradiction route}
The correctly acquired consequence of (1) is a polynomial, post-disclosure,
source-matched responder with exponent (4). The intended forbidden outcome
was an advance predictor. That final arrow is still \textbf{OPEN}:
\[
 A\in\mathrm P\text{ deciding }L
 \Longrightarrow B_A\text{ as in Theorem 1}
 \quad\not\Longrightarrow\quad
 \text{successful predictor of an unreleased block}.    \tag{7}
\]
The palette example falsifies automatic conversion, including when the
challenge genuinely changes the required witness. A successful continuation
of this route must prove a source-specific commitment-forcing theorem with
information access, uniformity and resources unchanged. No general mandatory
implementation compiler is required by every possible P-vs-NP argument;
it is this particular live-prediction route that needs its proposed transfer.

A direct non-P proof for $L$, or a model-matched resource contradiction,
remains allowed. Theorem 1 does not eliminate the arbitrary exponent $k$.
Repeated oracle calls that are polynomially bounded in number and size
preserve polynomial time, and fixed additive or multiplicative overhead does
not contradict (1). Hashes remain valuable for exact accounting and replay,
not an assumed source of cryptographic or information-theoretic hardness.

\section{Bounded evidence}
Four independent invocations, each with a 25-second guard, passed
\textbf{179,370 assertions}, with zero failures and zero timeouts. Their
recorded times total 1.449 seconds; the longest was 0.573 seconds rounded up.
No random seed was used. These are finite tests, not a polynomial SAT solver.
\begin{center}\small
\begin{tabular}{lrr}\toprule
Batch & Assertions & Seconds\\\midrule
All four-bit transcript relations &131,074&0.512\\
Source/target trees and prefixes &41,133&0.364\\
Active hash/pin histories and exact replay &7,007&0.572\\
Horn/Krom/affine/product controls &156&0.001\\\bottomrule
\end{tabular}\end{center}
The tree batch covers all 75 labelled graphs on one through four vertices,
1,776 source colourings, 20,826 palette repairs, 8,094 prefix instances,
224 independent target backtracking decisions and 185 reverse witnesses.
The target solver is distinct from the canonical forward lift. The finite
prefix oracle scans full source solution lists: its 746,401 examined rows
are explicitly recorded, not charged as constant-time decisions.

The hash batch covers 1,152 three-round fixtures with 2,620 events and 2,080
verified witnesses. Both replaceable single pins and retained multiple pins
are tested. It records 540 local NO\_EXTENSION events, 15,738 finite oracle
queries and 116,087 source-solution rows scanned. SHA-256 is used at widths
256, 2 and zero bits. The zero-bit case is a constant digest: exact-record
checks still work, while digest equality alone says nothing. All replay
checks receive the actual supplied sample tape. They do not authenticate
physical sampling from an untrusted event record.

Binary source controls include Horn, dual-Horn, 2-SAT, affine parity,
bounded-treewidth paths, independent products, unused-variable padding and
an inconsistent formula. No operation count for these solvers becomes a
minimum-work claim. COPY/replay and source-aware reconstruction are explicit;
no primitive trace or wall-clock benchmark is reported.

\section{Barrier and corpus audit}
\paragraph{Literature first.}
The selected obstacles were post-release response acquisition and the
static-to-causal quantifier transfer. Primary sources checked include Cook's
reducibility theorem, D\"ocker's exact occurrence language, Chandra--Kozen--
Stockmeyer alternation, Baker--Gill--Solovay relativization, Aaronson--Wigderson
algebrization, and the Mudigonda--Williams joint-resource bounds
\cite{cook,docker,cks,bgs,aw,mw}. Search-to-decision and causal games are direct
model-matched imports. State entropy and fixed proof-system size are not used
as unrestricted time measures.

\paragraph{Oracle and algebraic audit.}
The finite causal impossibility survives an oracle independent of the future
sample. Under an equality oracle $P^O=NP^O$, the ordinary prefix language
(3) is also solvable in polynomial oracle time, while advance prediction
still fails. No contradiction arises because (7)'s last arrow is unproved.
Ordinary $L$ is not silently called NP$^O$-complete; its original reduction
still suffices for ordinary (3). No nonrelativizing/nonalgebrizing hinge is
claimed. Natural Proofs is not invoked. Hash security, ETH and random-oracle
hardness are not assumed. The useful PQ transfer is polynomially verified
preimage/response search; cryptographic binding as an unconditional cost
premise is a \textbf{FAILED TRANSFER}.

\paragraph{CORPUS REVIEWED.}
Fresh reads include AGENTS.md, CODEX.md, the research skill and mode guide,
root main.tex through its complete bridge/synthesis region, actual main4
and relevant main5 driver text, the recursive tree survey (truncated), task
directory and TASK-004/005, AFFINE-RESIDUAL-REDIRECT1, COMPLETIONCUT1,
CONTGAME1 and the exact two-tree construction/encoding code. The current root
records the latest TWOTREEBRIDGE1 MAIN. Its source bridge code and both
witness directions were inspected and independently tested here. The latest
LSPV-3 directory commit points to VALSUBCUBEFW1; no missing scientific body
was reconstructed. MODELBRIDGE-0022's canonical path and original header
were confirmed present, so no duplicate recovery was attempted.
Local git failed DNS resolution. No full checkout, repository-wide preflight
or CI, exhaustive historical survey, or fresh original Drive-proof reread
is claimed. Unread task sections are not certified as reviewed. This is a
bounded source review, with the limitation preserved in the outbox.

\paragraph{DUPLICATION / CONTRADICTION AUDIT.}
CONTGAME1/ROUNDABSORB1 already contain the abstract local-PH versus global
continuation boundary. COMPLETIONCUT1 already separates future prediction
from post-release work and specifies actual sample replay. TWOTREEBRIDGE1
already supplies a conditional extension checker under P=NP. These are not
rebranded as new discoveries. This report makes their composition explicit
for an active hash-dependent verifier, supplies the arbitrary-$k$ ledger,
and tests source-matched palette substitution, independent target witnesses,
and all four-bit causal/offline relations. No literature-wide novelty claim.

\paragraph{EARLY22 / SINGLE22 acknowledgment.}
\textbf{Already integrated} at root main.tex, Historical interfaces preserved;
the original recovered proofs were not reaudited this run. Signed exact-$(2,2)$
compiler/minimality/Resolution claims remain distinct from clause-monotone
cubic-pair representations and affine-continuation limits. The present use of
$L_{\mathrm{mono22}}$ relies on its published NP-completeness, not an
unproved monotonicity property of the signed compiler. No proof projection
is promoted to a decider lower bound. Exact new location: Sections 1--4 here
and the proposed LIVEPREFIX1 paragraph inside LSPV-4 ownership. One next
check is the commitment-forcing arrow (7), not reconstruction of old proofs.

\paragraph{WHAT CHANGED RELATIVE TO ROOT MAIN.}
The bridge now distinguishes exact active-constraint response from metadata
replay, and supplies a source-level recolouring substitution with positive
and negative target controls. It adds no new restricted-proof lower bound
and does not change the all-algorithm OPEN status.

\paragraph{ADOPTED / REJECTED / RETIRED PRIOR COMPONENTS.}
Adopted: exact source/target maps, COMPLETIONCUT1, CONTGAME1, full sample
accounting and source-aware recomputation. Rejected: swapping offline and
causal quantifiers, treating a changed reference hash as a semantic failure,
or silently importing small-space bounds into a polynomial-workspace
execution. Automatic commitment forcing is \textbf{FALSIFIED/RETIRED}; an
explicit source-specific transfer remains \textbf{OPEN}.

\paragraph{STALE OR OVERSTATED CLAIMS FOUND.}
No valid older theorem is retracted. The statement ``putting the hash inside
the constraints prevents the polynomial selector'' is false under the stated
bounded local-verifier assumptions. The stronger claim for arbitrary global
continuation games is not made in either direction.

\paragraph{Publication discipline and next target.}
Exact code, four raw batches and this body are preserved before publication.
Canonical report and root success require separate readbacks; the publication
manifest records actual outcomes. Only the LSPV-4 ownership region may change.
No schedule or team membership is modified. The next target is a concrete
commitment-forcing transformation on the two-tree source that preserves the
information boundary and every cost, rather than requiring advance knowledge
by definition. P versus NP remains \textbf{OPEN}.

\begin{thebibliography}{9}
\bibitem{docker} J. D\"ocker. Monotone 3-Sat-(2,2) is NP-complete.
\url{https://arxiv.org/abs/1912.08032}.
\bibitem{cook} S. A. Cook. The complexity of theorem-proving procedures.
STOC 1971, 151--158. \url{https://doi.org/10.1145/800157.805047}.
\bibitem{cks} A. K. Chandra, D. C. Kozen and L. J. Stockmeyer.
Alternation. JACM 28(1), 114--133, 1981.
\url{https://doi.org/10.1145/322234.322243}.
\bibitem{bgs} T. Baker, J. Gill and R. Solovay. Relativizations of the
P=?NP Question. SIAM J. Comput. 4(4), 431--442, 1975.
\url{https://doi.org/10.1137/0204037}.
\bibitem{aw} S. Aaronson and A. Wigderson. Algebrization: A New Barrier
in Complexity Theory. ECCC TR08-005, 2008.
\url{https://eccc.weizmann.ac.il/eccc-reports/2008/TR08-005/index.html}.
\bibitem{mw} A. S. Mudigonda and R. R. Williams. Time-Space Lower Bounds
for Simulating Proof Systems with Quantum and Randomized Verifiers.
\url{https://arxiv.org/abs/2012.00330}.
\end{thebibliography}
\end{document}
% END_OF_VERIFIED_REPORT::LSPV-4::LIVEPREFIX1
